% =====================================================================
% PORCE — Operacionalizacion regulatoria EASA y trazabilidad (paste-ready)
% Preambulo requerido (ya en Main_formato_ieee.tex): booktabs, siunitx, tabularx.
% Objetivo: que la motivacion EASA este reflejada EN EL ALGORITMO, no solo en la
% intro, y sea verificable contra el codigo. Acompana a la implementacion
% R_s(clase) (flight_controller._obs_safety_radius_m / porce_manager.plan_route).
% =====================================================================

\subsection{From Regulatory Motivation to Algorithmic Mechanism}
\label{subsec:reg-operationalization}
A recurring weakness of inspection planners is \emph{obstacle homogeneity}: a
tower, an animal, and an uninvolved person are inflated by the same margin, so
the European distinction between an asset and a protected third party never
reaches the planner. The proposed method removes that homogeneity at the one
place where it changes behaviour --- the obstacle inflation radius. Perception
classes are canonicalized into an operational family (\texttt{person},
\texttt{bicycle}, \texttt{biker}~$\rightarrow$ the uninvolved-person family
\texttt{bike}; \texttt{cow}; \texttt{tower}), and each family is assigned a
distinct safety radius $R_s(c)$ that the bounded A$^\star$ uses to carve the
local occupancy grid (Sec.~\ref{sec:formulation}, Eq.~\eqref{eq:grb}). The
uninvolved-person radius is not hand-tuned: it is the SORA \emph{Ground Risk
Buffer} under the conservative 1:1 rule, i.e. a horizontal protection at least
equal to the operating height, clipped to an operational range. Towers receive
the smallest clearance because inspection deliberately flies close to the asset;
unrecognized classes default to the person radius, so the failure mode is
conservative. Table~\ref{tab:class-radius} lists the resulting values and
Table~\ref{tab:traceability} maps each regulatory concept to the concrete
mechanism and the repository artifact that implements it.

\begin{table}[t]
    \centering
    \small
    \begin{tabular}{@{}l l l@{}}
        \toprule
        \textbf{Class (canonical)} & \textbf{$R_s$} & \textbf{Regulatory basis} \\
        \midrule
        Uninvolved person (\texttt{bike}) & $\mathrm{clip}(h,15,40)\,\si{\meter}$ & SORA GRB, 1:1 rule \\
        Animal (\texttt{cow}) & \SI{12}{\meter} & Physical obstacle (fixed) \\
        Asset (\texttt{tower}) & \SI{8}{\meter} & Inspection target (fixed) \\
        Unknown & as person & Conservative default \\
        \bottomrule
    \end{tabular}
    \caption{Class-dependent safety radius $R_s(c)$. At the validated cruise
    height $h\approx\SI{23}{\meter}$ the person radius evaluates to
    \SI{23}{\meter}, strictly larger than the animal and asset clearances.}
    \label{tab:class-radius}
\end{table}

\begin{table*}[t]
    \centering
    \small
    \begin{tabular}{@{}p{0.27\textwidth} p{0.40\textwidth} p{0.27\textwidth}@{}}
        \toprule
        \textbf{EASA / SORA concept} & \textbf{System mechanism (what the algorithm actually does)} & \textbf{Repository artifact} \\
        \midrule
        Uninvolved person (Reg. (EU) 2019/947; SORA 2.5) &
        Civilian-like detections canonicalized into the \texttt{bike} family and given the largest, dedicated protection radius; unknown classes default to it &
        \texttt{flight\_controller.py}: \texttt{\_canonical\_obs\_class\_name}, \texttt{\_obs\_safety\_radius\_m} \\

        Ground Risk Buffer / 1:1 rule (SORA Annex F) &
        Person radius set to the height-dependent buffer $R_s=\mathrm{clip}(\kappa h,R_{\min},R_{\max})$, $\kappa=1$; grows with altitude like the regulatory buffer &
        \texttt{constants.py} (\texttt{SAFETY\_GRB\_RATIO}, floor/max); \texttt{\_obs\_safety\_radius\_m} \\

        Obstacle non-homogeneity (asset $\neq$ person) &
        Per-obstacle inflation: each grid seed inflated by its own $n_s(o)=\lceil R_s(c_o)/\delta\rceil$, so person $>$ animal $>$ asset &
        \texttt{porce\_manager.py}: \texttt{plan\_route}; \texttt{planner\_obstacle\_subset} tags \texttt{safety\_m} \\

        M1(C) Ground observation mitigation (avoid overflight of detected people) &
        Onboard detector observes the corridor; a detected person triggers a bounded local bypass instead of overflight; the maneuver is logged &
        \texttt{vision\_system.py} (YOLO + geo); evasion branch + \texttt{events.jsonl} \\

        No-overflight objective &
        A$^\star$ removes from the candidate corridor every cell inside a protected region, so the executed path does not pass within $R_s(\text{person})$ of an uninvolved person &
        \texttt{porce\_manager.py}: \texttt{plan\_route} occupancy + search \\

        Auditability / bounded behaviour (EASA AI Roadmap 2.0) &
        Deterministic graph search (no black-box policy) with a per-decision zero-trust audit trail usable as compliance evidence &
        \texttt{flight\_controller.py}: \texttt{decision\_snapshot}, \texttt{evasion\_route\_generated} \\
        \bottomrule
    \end{tabular}
    \caption{Regulatory traceability: each EASA/SORA concept is mapped to the concrete algorithmic mechanism and the repository artifact that implements it. This is the differentiator of the method --- it brings protected-ground logic, controller integration, and auditability into one loop rather than treating them as separate planning, perception, and compliance layers.}
    \label{tab:traceability}
\end{table*}

\subsection{No-Overflight Exposure Metric}
\label{subsec:overflight-metric}
To make the regulatory objective measurable rather than rhetorical, we define an
\emph{overflight-exposure event} at control step $t$ as any uninvolved-person
track whose horizontal distance to the vehicle falls below its protection radius,
\begin{equation}
E(t)=\Big|\{\,o\in\mathcal{O}_t : c_o=\texttt{person}\ \wedge\ \lVert p_t-p_{o}\rVert_2 < R_s(c_o)\,\}\Big|,
\end{equation}
and the mission exposure as $E_{\text{tot}}=\sum_t \mathbf 1[E(t)>0]$. The
operational claim of the method is that enabling class-dependent replanning drives
$E_{\text{tot}}$ toward zero relative to the asset-agnostic baseline, at the small
path/time overhead reported in the ablation. $E_{\text{tot}}$ is reconstructible
from the existing trajectory and track logs, so it can be reported without new
experiments beyond re-running the ablation under the class-dependent radii.
